\subsection{Conclusiones}
Los resultados obtenidos a lo largo del trabajo permiten afirmar 
que la compresión sin pérdida representa una alternativa válida 
para reducir la huella de memoria en simuladores cuánticos de estado 
completo sin alterar la exactitud numérica de la simulación. 
Su incorporación dentro del simulador mostró que este tipo de estrategias 
puede aplicarse de manera efectiva sin abandonar la representación 
completa del estado, lo que preserva el valor de estos simuladores 
como referencia precisa para el análisis y la validación de algoritmos cuánticos.

Al mismo tiempo, el comportamiento observado durante la evaluación 
dejó claro que el beneficio de la compresión no es uniforme. 
Su utilidad depende de la estructura del estado cuántico y del 
tipo de circuito que se ejecuta, de modo que existen escenarios 
en los que la reducción de memoria resulta favorable y otros en 
los que esa ventaja es más limitada. Más que una solución general 
para cualquier simulación, la compresión sin pérdida se perfila 
así como una estrategia especialmente adecuada cuando el estado 
conserva regularidades que pueden aprovecharse durante el almacenamiento.

En este sentido, uno de los aportes centrales del trabajo fue 
mostrar que el problema de memoria no se resuelve únicamente 
con la elección de una biblioteca de compresión, sino con una 
integración coherente entre particionamiento del vector, 
acceso selectivo y reutilización temporal de datos. 
La estrategia desarrollada permitió materializar esa idea en una 
implementación concreta y comprobar que la compresión puede 
incorporarse como parte de la gestión de memoria del simulador, 
y no solo como un mecanismo externo de almacenamiento.

En conjunto, el trabajo demuestra que la compresión sin pérdida 
puede ampliar la capacidad práctica de simulación exacta cuando 
la memoria constituye la restricción principal y el comportamiento 
del estado ofrece condiciones favorables para ello. Su principal 
alcance está en aportar una alternativa que conserva igualdad 
exacta respecto a la ejecución no comprimida y que, por esa misma razón, 
resulta útil en contextos donde no es aceptable sustituir precisión por aproximación.